% ============================================================================
% INTRODUCTION
% ============================================================================
\section{Introduction}

Most computational systems enforce a strict separation between three
concerns: memory stores data, programs specify how data is transformed,
and learning mechanisms adjust parameters to improve the transformations
over time.  This division is foundational in both classical software
engineering and modern machine learning, where data resides in tensors,
programs are defined by network architectures, and learning is driven by
gradient descent over differentiable loss surfaces.

This work explores an alternative formulation in which these distinctions
are collapsed.  We describe a four-layer architecture in which data,
programs, execution state, and learning coexist within a unified
substrate:

\begin{enumerate}[nosep]
  \item \textbf{Values} (Layer~1): fixed-width 1024-byte binary frames
        carrying token data, a locality-sensitive affinity fingerprint,
        an embedded program, Boolean and geometric execution lanes, and
        structural links---simultaneously data, code, and execution
        context.
  \item \textbf{MarkovTrie} (Layer~2): an adaptive probabilistic suffix
        trie that learns from every observation via surprise-modulated
        plasticity, requiring no training epochs, loss functions, or
        weight matrices.
  \item \textbf{Kadabra DHT} (Layer~3): a Kademlia-style distributed
        hash table that routes Values to MarkovTries by affinity
        similarity, creating emergent specialization through
        content-addressable clustering.
  \item \textbf{The Field} (Layer~4): a collective attention mechanism
        that detects emergent eigenmodes from gossip-propagated
        adaptive-state digests and projects asymmetric pressure onto
        each trie---amplifying learning in the dominant mode and
        suppressing it elsewhere.
\end{enumerate}

Computation at Layer~1 is a two-lane in-frame ALU.  The low opcode
nibble applies the sixteen two-input Boolean truth tables across
designated regions of one or more frames.  The high opcode nibble
dispatches Projective Geometric Algebra (PGA) operations---composition,
sandwich application, and reversal---over 64-byte multivectors stored in
the Context, Gradient, and Signals regions.  Because programs and
operands are stored within frames, there is no boundary between the
instruction stream and the data it operates on; a frame may serve as
passive storage, an executable program, a geometric motor, or all three.

Learning at Layer~2 is driven entirely by surprisal.  The trie's
learning rate, decay rate, prune threshold, classification context,
temperature, beam width, interpolation weights, episodic blend, and
unsupervised threshold are all derived online from signals the trie
already computes---no hyperparameter search is required.

Routing at Layer~3 exploits a SimHash projection of the token content
into the affinity region: five independent locality-sensitive hash
projections using prime strides populate a 257-bit address so that
Values with similar content are routed to the same network node, where a
single MarkovTrie accumulates expertise in that content domain.

Attention at Layer~4 emerges from the interaction of gossip-propagated
digests.  Each node broadcasts a compact summary of its surprisal,
growth rate, finite-field phase, and affinity vector.  The Field detects
clusters of nodes whose affinity vectors clear an adaptive Jaccard
coupling threshold, then uses surprisal-velocity phase coherence and the
global finite-field phase to modulate each node's decay and learning
multipliers toward the dominant pattern.

The architecture does not employ gradient-based optimization or learned
parameters.  Its fast path remains dense Boolean logic---XOR distance,
population count, and truth-table evaluation---but the current runtime
also uses localized floating-point geometry for PGA motors, Orthogonal
Procrustes alignment, and episodic drift correction.  A multi-substrate
backend routes both Boolean and geometric frame operations to whichever
compute hardware is available---CPU, NVIDIA CUDA, or Apple Metal---
through a unified kernel interface.

We evaluate the system across five experimental domains: text
classification, multi-language code generation, CIFAR-10 image
reconstruction, logical reasoning on the bAbI benchmark, and several
text generation tasks including compositional assembly, out-of-corpus
retrieval, and prose chaining.  The experiments are defined
programmatically, executed within the substrate, and compiled into the
present document automatically.

The contributions of this paper are as follows.  First, we define the
four-layer architecture: the Value frame, the MarkovTrie, the Kadabra
DHT, and the Field---and show how they compose into a functioning
runtime that learns without gradients.  Second, we provide the
mathematical foundations: SimHash collision bounds, affinity coupling
for eigenmode detection, surprise-modulated plasticity, PGA transition
operators, and Orthogonal Procrustes alignment for cross-domain
projection.  Third, we describe the native kernel path that preserves
the Boolean filter while adding CPU, CUDA, and Metal execution for the
geometric lane.  Fourth, we report empirical results across modalities,
providing a baseline characterization of what this class of
architecture can and cannot achieve.  Fifth, we demonstrate a
self-documenting experimental pipeline in which the system produces
both results and the corresponding research artifact.

This work does not aim to match the performance of large parameterized
models on standard benchmarks.  Its purpose is to explore a different
point in the design space: one where retrieval and filtering reduce to
deterministic, hardware-native bitwise operations; where continuous
geometry is introduced only where it supplies an explicit spatial
operator; where learning is local and surprisal-driven; and where
attention emerges from collective dynamics rather than from learned
projection matrices.
